\section*{Chapitre \ref{ch:in:asset-managers-pensions-sovereign-funds-and-insurers} --- Sociétés de gestion, régimes de retraite, fonds souverains et assureurs}

\begin{solution}{exo:in:asset-managers-pensions-sovereign-funds-and-insurers:1}
$1.1300/11.7173=0.09644$ dollars par couronne~: 21\,268~milliards de NOK valent 2\,051~milliards de dollars~; par salarié, $2\,051/678=\$3.03$
milliards de dollars.
\end{solution}

\begin{solution}{exo:in:asset-managers-pensions-sovereign-funds-and-insurers:2}
$1\,756/1\,062\,000=16.5$ points de base.
\end{solution}

\begin{solution}{exo:in:asset-managers-pensions-sovereign-funds-and-insurers:3}
Le sponsor~: la prestation est fixe, donc un déficit des actifs retombe sur lui. Le travail du quant relève de la gestion adossée au
passif (LDI)~: valorisation des engagements, appariement de la duration et de l'inflation, couverture et mesure du risque sur le ratio de financement.
\end{solution}

\begin{solution}{exo:in:asset-managers-pensions-sovereign-funds-and-insurers:4}
Frais externes $(1\,976+2\,758)/793\,300=59.7$ points de base~; charges d'exploitation $1\,757/793\,300=22.1$. Le 23.1 publié
repose sur la définition et le dénominateur propres au fonds (actifs nets moyens sur l'année)~; un dénominateur de fin d'année est
plus élevé lors d'une année de croissance.
\end{solution}

\begin{solution}{exo:in:asset-managers-pensions-sovereign-funds-and-insurers:5}
Coût $20\times0.352+1.5=8.54$ millions de dollars~; seuil de rentabilité $8.54/0.0025=3.42$ milliards de dollars.
\end{solution}

\begin{solution}{exo:in:asset-managers-pensions-sovereign-funds-and-insurers:6}
L'équipe coûte 5.02~millions de dollars~: $5.02/2\,000\times10^4=25.1$ points de base.
\end{solution}

\begin{solution}{exo:in:asset-managers-pensions-sovereign-funds-and-insurers:7}
La commission attendue vaut $20+0.2\times100=40$ points de base~: seuil de rentabilité $5.02/0.004=1.255$ milliards de dollars. Cela ne tient que si
l'équipe interne dégage aussi 1~\% de rendement excédentaire~; sinon, il faut comparer les coûts nets des rendements.
\end{solution}

\begin{solution}{exo:in:asset-managers-pensions-sovereign-funds-and-insurers:8}
Les deux fonds ont des modèles économiques différents~: l'un détient des actifs cotés proches des indices de référence, gérés en interne~; l'autre
achète des stratégies privées et actives dont les commissions financent des rendements excédentaires attendus. Le coût n'est que la
moitié de la comparaison~; l'autre moitié est le rendement net, et la composition des actifs diffère.
\end{solution}

\begin{solution}{pb:in:asset-managers-pensions-sovereign-funds-and-insurers:1}
\begin{enumerate}
\item Une institution qui investit pour le compte de bénéficiaires ou de l'État~; un \omterm{def:in:asset-managers-pensions-sovereign-funds-and-insurers:swf}{fonds souverain}, détenu par l'État, à des fins macroéconomiques, investissant en actifs
étrangers~; un \omterm{def:in:asset-managers-pensions-sovereign-funds-and-insurers:swf}{régime de retraite à prestations définies} promettant une prestation fixe à la retraite.
\item Gestion par les équipes propres du propriétaire~; 95~\%.
\item Allocation d'actifs, risque, appariement au passif, actions factorielles et quantitatives, gestion indicielle et mise en œuvre.
\item 3.8 et 23.1 points de base~; l'un inclut les frais des gérants externes et l'autre les exclut, avec des
dénominateurs différents.
\item Personnel, conservation, informatique et données, recherche et juridique, autres coûts et coûts alloués, ainsi que les frais de base et
commissions de performance des gérants externes.
\item 3.5 et 81.8 points de base.
\item 2.3 points de base sur les actifs gérés en interne~; 27.7 sur les mandats externes.
\item $59.7/81.8=73~\%$.
\item 3.65~millions de NOK~; 352\,000~dollars.
\item 3.0~milliards de dollars.
\item Coût interne rapporté au taux de frais.
\item 5.02~millions de dollars par an.
\item 5.02, 2.51 et 1.255~milliards de dollars.
\item $15/2.51=6.0$ fois.
\item 30~millions de dollars par an.
\item 5.02~milliards de dollars, 2.51~milliards de dollars et 1.255~milliards de dollars.
\item Le même rendement net de coûts que celui du gérant externe.
\item Ce que paient les coûts~: la composition des actifs, les stratégies et les rendements nets.
\item Elle crée des emplois de gestion de portefeuille, de trading et de recherche~; elle réduit le travail de sélection des gérants et de suivi du
mandat.
\item À 20 points de base, le mandat vaut six fois le seuil de rentabilité, donc la \omterm{def:in:asset-managers-pensions-sovereign-funds-and-insurers:internal}{gestion interne} est bien moins chère si l'équipe peut
égaler le rendement net du gérant~; la décision repose sur cette capacité, et non sur le coût.
\end{enumerate}
\end{solution}

\begin{solution}{iq:in:asset-managers-pensions-sovereign-funds-and-insurers:1}
Un propriétaire répond devant ses bénéficiaires ou un gouvernement et choisit ce qu'il gère et à qui il confie la gestion~; un gérant touche des frais
pour gérer des mandats. Chez un propriétaire, un quant travaille sur l'allocation, les passifs, le risque et la sélection des gérants~; chez un gérant, sur
la production du rendement fixé par le mandat.
\iqlookfor{principal contre agent, et le travail qui en découle.}
\end{solution}

\begin{solution}{iq:in:asset-managers-pensions-sovereign-funds-and-insurers:2}
Les passifs augmentent d'environ 18~\% pour une baisse de 1~\% des taux, et les obligations d'environ 6~\%, donc le ratio de financement baisse. Allonger la
duration des actifs ou ajouter des swaps de taux (recevoir le taux fixe) pour couvrir une partie de l'écart.
\iqlookfor{décalage de duration et couverture.}
\end{solution}

\begin{solution}{iq:in:asset-managers-pensions-sovereign-funds-and-insurers:3}
L'excédent net pour le propriétaire vaut $x-50-0.2x$ points de base, contre $-5$ pour le portefeuille indiciel~: on retient le gérant si $0.8x-50>-5$, soit
$x>56.25$ points de base par an, avant tout ajustement du risque.
\iqlookfor{frais sur le rendement excédentaire, et une inégalité simple.}
\end{solution}

\begin{solution}{iq:in:asset-managers-pensions-sovereign-funds-and-insurers:4}
L'allocation fixe l'essentiel du risque du portefeuille, donc la dispersion de ses rendements~; la sélection au sein des classes d'actifs ajoute
un montant plus faible et moins certain. Les meilleurs quants devraient éclairer l'allocation et le risque, et pas seulement la sélection de titres.
\iqlookfor{la décomposition du risque indique où l'effort paie.}
\end{solution}

\begin{solution}{iq:in:asset-managers-pensions-sovereign-funds-and-insurers:5}
Courbes, duration et convexité, produits dérivés de taux et couverture, XVA et transfert de collatéral~; à acquérir~: modèles de passifs,
mortalité et rachats, et règles de capital de solvabilité.
\iqlookfor{compétences transférables en taux et nouveau volet passif.}
\end{solution}

\begin{solution}{iq:in:asset-managers-pensions-sovereign-funds-and-insurers:6}
Régresser les rendements excédentaires mensuels de l'équipe, nets de ses coûts, sur l'indice de référence (et les facteurs), estimer l'alpha avec une
erreur robuste à l'hétéroscédasticité et à l'autocorrélation, et se demander si dix ans donnent assez de puissance~: avec une tracking error de
2~\% et un alpha vrai de 0.5~\%, la statistique t après dix ans vaut environ 0.8.
\iqlookfor{net de coûts, erreurs robustes et puissance.}
\end{solution}
